\begin{thebibliography}{99}
\raggedright
\bibitem{timetravel} Lianhui Qin et al. Counterfactual Story Reasoning and Generation. EMNLP-IJCNLP, 2019.
\url{https://aclanthology.org/D19-1509/}.
\bibitem{pasta} Sayontan Ghosh et al. PASTA: A Dataset for Modeling PArticipant STAtes in Narratives. TACL, 2023.
\url{https://aclanthology.org/2023.tacl-1.73/}.
\bibitem{dto} Amelie Girard and Massimo Piccardi. DTO: a Differentiable Training Objective for Effective Counterfactual Story Rewriting. arXiv preprint, 2026.
\url{https://arxiv.org/abs/2605.24885}.
\bibitem{cfi} Basar Demir et al. Aligning Language Models Towards Counterfactual Invariance. arXiv preprint, 2026.
\url{https://arxiv.org/abs/2610.22942}.
\bibitem{csr} Sanjeda Akter, Ibne Farabi Shihab and Anuj Sharma. Counterfactual Sensitivity Regularization for Trace-Sensitive Reasoning in Large Language Models. Findings of EMNLP, 2026.
\url{https://aclanthology.org/2026.findings-emnlp.539/}.
\bibitem{checklist} Marco Tulio Ribeiro et al. Beyond Accuracy: Behavioral Testing of NLP Models with CheckList. ACL, 2020.
\url{https://aclanthology.org/2020.acl-main.442/}.
\bibitem{openmeva} Jian Guan et al. OpenMEVA: A Benchmark for Evaluating Open-ended Story Generation Metrics. ACL, 2021.
\url{https://aclanthology.org/2021.acl-long.500/}.
\bibitem{mingru} Leo Feng et al. Were RNNs All We Needed? arXiv version, 2024.
\url{https://arxiv.org/abs/2410.01201}.
\bibitem{tiny} Ronen Eldan and Yuanzhi Li. TinyStories: How Small Can Language Models Be and Still Speak Coherent English? arXiv preprint, 2023.
\url{https://arxiv.org/abs/2305.07759}.
\bibitem{fairytaleqa} Ying Xu et al. Fantastic Questions and Where to Find Them: FairytaleQA---An Authentic Dataset for Narrative Comprehension. ACL, 2022.
\url{https://aclanthology.org/2022.acl-long.34/}.
\end{thebibliography}
